\documentclass[conference]{IEEEtran}
\IEEEoverridecommandlockouts
\usepackage{cite}
\usepackage{amsmath,amssymb,amsfonts}
\usepackage{algorithmic}
\usepackage{graphicx}
\usepackage{textcomp}
\usepackage{xcolor}
\usepackage{booktabs}
\usepackage{hyperref}

\begin{document}

\title{Evaluación de Robustez y Ruido en Detección de Objetos en el Mundo Real: Defensa Contra la Degradación de Sensores}

\author{\IEEEauthorblockN{William Steve Rodriguez Villamizar}
\IEEEauthorblockA{\textit{AI Leader \& Solutions Architect} \\
\textit{wisrovi-suit} \\
Badajoz, Spain \\
wisrovi.rodriguez@gmail.com}
}

\maketitle

\begin{abstract}
Los modelos de detección de objetos desplegados en entornos del mundo real a menudo encuentran degradación de sensores, como ruido gaussiano, desenfoque de movimiento y artefactos inducidos por el clima. Este documento evalúa la robustez de la arquitectura YOLOv8 contra el ruido de sensores sintéticos, cuantificando la degradación de la Precisión Media (mAP) bajo diferentes intensidades de ruido. A través de una simulación de micro-benchmark controlada en el conjunto de datos COCO, demostramos que las arquitecturas estándar sufren una caída precipitada en la sensibilidad (recall) cuando se exponen a ruido gaussiano de alta varianza. Proponemos un marco medido para la evaluación de ruido para apoyar despliegues confiables en escenarios de edge-computing, proporcionando evidencia para su despliegue.
\end{abstract}

\begin{IEEEkeywords}
YOLO, Detección de Objetos, Evaluación de Ruido, Degradación de Sensores, Robustez, MLOps
\end{IEEEkeywords}

\section{Introducción}
Los despliegues en el mundo real de modelos de detección de objetos a menudo sufren de cambios de dominio causados por la degradación física de los sensores \cite{hendrycks2019robustness}. Mientras que modelos como YOLOv8 \cite{jocher2023ultralytics} sobresalen en benchmarks impecables como COCO \cite{lin2014microsoft}, su rendimiento puede degradarse drásticamente bajo condiciones de ruido gaussiano o baja iluminación. Este estudio cuantifica esta degradación a través de un marco empírico sistemático.

\section{Metodología}
Simulamos la degradación de los sensores inyectando ruido gaussiano ($\mu=0$, $\sigma \in [0.1, 0.5]$) en el conjunto de validación COCO128. Luego evaluamos el rendimiento de YOLOv8 a través de estos regímenes de ruido.

\section{Resultados Experimentales}
Nuestra simulación de micro-benchmark controlada indica que el mAP50 se degrada de 0.605 en datos impecables a 0.421 bajo ruido moderado ($\sigma=0.3$) y 0.153 bajo ruido severo ($\sigma=0.5$).

\section{Conclusión}
Nuestro marco medido proporciona evidencia de que evaluar explícitamente modelos contra ruido de sensores es necesario para apoyar despliegues confiables en entornos industriales.

\section*{Disponibilidad de Datos y Código}
Los scripts y sus resultados empíricos estrictamente ejecutados en CSV se publican en la carpeta \texttt{evidencias/} de este documento. Este ecosistema opera bajo un Modelo de Licencia Dual (PolyForm Noncommercial / AGPLv3). 

\section*{Agradecimientos}
Este trabajo fue apoyado por wisrovi-suit.

\bibliographystyle{IEEEtran}
\bibliography{references}

\end{document}
